\section{An explicit witness: the separation on a rational datum}\label{sec:witness}

\subsection{The witness datum}\label{sec:datum}

Two architectures -- extraction ($q=0$) and regeneration ($q=1$) -- share one review interval $[0,1]$. The phase state is $(q,x,s_1,s_2)$: a physical reserve stock $x$ (a bridge stock, for instance an adjustment fund) and two typed floors $s_1$ and $s_2$ normalised at zero (on the benchmark of Section~\ref{sec:benchmark}: a biomass margin above a limit reference point and a fleet income margin). The transition-safe set is $S_0=\{x\ge0,\ s_1\ge0,\ s_2\ge0\}$ and the destination set is $G=\{(1,x,s):x\ge0,\ s\ge0\}$, with destination maintainability witnessed by a hold policy. The destination reset adds the gain $e=(\tfrac14,\tfrac14)$ to both floors; $e$ is strictly positive, so successors are interior to $G$ and no inequality below binds on its magnitude. The rescue cost is $c=1$. Initial states are $X_0=\{(0,x,s):x\ge0,\ s\ge0\}$.

Four plans are available (Table~\ref{tab:actions}), written as piecewise-linear maps on the breakpoints $\{0,\tfrac12,1\}$ and monotone on each piece, so that every tube is the exact visited set. The disturbance class is \emph{action-indexed}: the adverse disturbance triggers the worst-case dip, of fixed depth $2$, in the active plan's characteristic coordinate ($s_1$ for \act{fast}, $s_2$ for \act{slow}); on the benchmark the depth $2$ is a quota pulse of $3/2$ plus an environmental shock of $1/2$. The worst-case tube of \act{fast} is therefore $[s_1-2,s_1]$ in $s_1$ with $s_2$ and $x$ constant, that of \act{slow} is symmetric, and that of \act{staged} is $[x-1,x]$ in $x$ with both floors nondecreasing. A shared-shock class, in which the environmental component strikes every plan, is treated in Section~\ref{sec:commonshock}.

\begin{table}[!htbp]
\centering\small
\caption{The plan menu of the witness datum from an initial state $(0,x,s)$, $x\ge0$, $s\ge0$. Trajectories are worst-case; $e=(\tfrac14,\tfrac14)$.}
\label{tab:actions}
\begin{tabularx}{\textwidth}{@{}lXl@{}}
\toprule
Plan & Within-interval trajectory (worst case) & Successor \\
\midrule
\act{no-switch} & state constant & $\{(0,x,s)\}$ -- misses $G$ \\
\act{fast} & $s_1(t)=s_1-4t$ on $[0,\tfrac12]$, $s_1-4(1-t)$ on $[\tfrac12,1]$; $s_2,x$ constant & $\{(1,x,s+e)\}$ \\
\act{slow} & $s_2(t)=s_2-4t$ on $[0,\tfrac12]$, $s_2-4(1-t)$ on $[\tfrac12,1]$; $s_1,x$ constant & $\{(1,x,s+e)\}$ \\
\act{staged} & $x(t)=x-t$ on $[0,1]$; floors grow linearly to $s+e$ & $\{(1,x-1,s+e)\}$ \\
\bottomrule
\end{tabularx}
\end{table}

The datum is deliberately spare. Its margin-class plans are \act{fast} and \act{slow}, with $\mathcal D=\{(2,0),(0,2)\}$, so Theorem~\ref{thm:menu} predicts the gap $\{s_1+s_2\ge2\}\setminus(\{s_1\ge2\}\cup\{s_2\ge2\})$ among them; \act{staged} lies outside the margin class, because its admissibility is governed by the reserve rather than by the floors, and it is what splits the predicted gap into a rescuable and an impossible part.

\subsection{The witnessed separation}

\begin{theorem}[witnessed separation]\label{thm:witness}
On the witness datum, over $X_0$:
\begin{enumerate}[label=\textup{(\arabic*)},leftmargin=2em]
\item $\Vtyp=\{x\ge1\}\cup\{s_1\ge2\}\cup\{s_2\ge2\}$.
\item $\Vweak=\bigcap_{w\in\Wp}\mathcal V_w=\{x\ge1\}\cup\{s_1+s_2\ge2\}$.
\item $\Vphys=\Vend=X_0$.
\end{enumerate}
Define the \emph{discrepancy region} $\Q=\{s_1<2,\ s_2<2,\ s_1+s_2\ge2\}$ ($x$ unrestricted), the \emph{rescue set} $R=\Q\cap\{x\ge1\}$ and the \emph{impossibility region} $I=\Q\cap\{x<1\}$. Then:
\begin{enumerate}[label=\textup{(\arabic*)},leftmargin=2em,resume]
\item $\FP=\Vweak\setminus\Vtyp=I$. The acceptance gap is exactly the impossibility region; it is nonempty, with relative interior $\{0<x<1,\ 0<s_1<2,\ 0<s_2<2,\ s_1+s_2>2\}$ in the $q=0$ slice. The rescue set $R$ lies in $\Vtyp$, not in the gap.
\item Both hierarchy inclusions are strict: every point of $I$ lies in $\Vweak\setminus\Vtyp$, and $(x,s_1,s_2)=(\tfrac12,\tfrac1{10},\tfrac1{10})$ lies in $\Vphys\setminus\Vweak$.
\item \textup{(Per-weight plan disagreement.)} On the relative interior of $\Q$ ($0<s_2<2$), with $r=w_2/w_1$, \act{fast} is licensed exactly for $r\ge\rho_1$ and \act{slow} exactly for $r\le\rho_2$, where
\[
\rho_1=\frac{2-s_1}{s_2},\qquad \rho_2=\frac{s_1}{2-s_2},\qquad \rho_2\ge\rho_1\iff s_1+s_2\ge2 .
\]
Assessors with $r>\rho_2$ license \act{fast} only, those with $r<\rho_1$ license \act{slow} only, intermediate assessors license both. On $I$ no single plan serves every weight: $\bigcap_wE_w(z)=\Etyp(z)=\varnothing$. On $R$, \act{staged} is typed-admissible and serves every weight.
\item \textup{(The rescue split.)} $R$ is typed-transformable, witnessed by \act{staged}: bridging at cost $c=1$ keeps both floors intact and lands in $G$. $I$ is aggregate-feasible for every cone weight yet admits no typed-admissible plan, with four exhibited violations: \act{fast} breaches the $s_1$ floor under its adverse disturbance, \act{slow} the $s_2$ floor, \act{staged} drives $x$ negative, and \act{no-switch} misses $G$.
\end{enumerate}
\end{theorem}

The proof is in \ref{app:witness}; the formulas for $\rho_1,\rho_2$ require $0<s_2<2$, and on the boundary faces the classification follows from the closed-set identities (1)--(2). Table~\ref{tab:regions} records the two verdicts region by region, and Figures~\ref{fig:triangle} and \ref{fig:paths} display them.

\paragraph{The witness point} At $(x,s_1,s_2)=(\tfrac12,\tfrac65,\tfrac65)$ the typed verdict is rejection: \act{staged}'s reserve tube reaches $\tfrac12-1<0$, \act{fast}'s worst-case $s_1$ dips to $\tfrac65-2=-\tfrac45$, \act{slow}'s $s_2$ likewise, and \act{no-switch} misses $G$. Yet $s_1+s_2=\tfrac{12}5\ge2$, so the state is compensatorily accepted, with thresholds $\rho_1=\tfrac23$ and $\rho_2=\tfrac32$: $r=2$ licenses \act{fast} only, $r=\tfrac12$ licenses \act{slow} only, $r=1$ licenses both. Every weighting licenses a plan; no plan is licensed by every weighting. With the same floors and $x=1$ the state lies in $R$: \act{staged}'s tube is $[0,1]$, the floors grow, and the successor $(1,0,s+e)\in G$.

\paragraph{Reading} The compensatory assessment's binding condition is a total budget, $s_1+s_2\ge2$; the noncompensatory assessment requires either one floor to survive its own worst-case dip ($s_i\ge2$) or the reserve to cover the rescue cost ($x\ge1$). Between the two conditions the compensatory doctrine certifies -- per weight, with weight-dependent plans -- transitions the noncompensatory doctrine rejects, and the aggregate reading cannot tell $I$ from $R$: the two differ only in $x$, which the index does not see. Part (7) is a \emph{negative certificate} in the sense of a rejection with an exhibited violated constraint for each plan, exhausting the menu; it is a complete verdict relative to the menu, not an inconclusive one.

\begin{table}[!htbp]
\centering\small
\caption{Region dictionary on the witness datum (Theorem~\ref{thm:witness}).}
\label{tab:regions}
\begin{tabularx}{\textwidth}{@{}>{\raggedright\arraybackslash}p{0.30\textwidth}>{\raggedright\arraybackslash}p{0.17\textwidth}>{\raggedright\arraybackslash}p{0.14\textwidth}X@{}}
\toprule
Region & Per-weight verdict & Common-plan verdict & Reason \\
\midrule
$x\ge1$, all $s$ & accept & accept & \act{staged} serves every weight \\
$s_1\ge2$ or $s_2\ge2$ & accept & accept & \act{fast}/\act{slow} survives its own dip \\
$R$: $x\ge1$, $s_1<2$, $s_2<2$, $s_1+s_2\ge2$ & accept & accept & the reserve finances \act{staged} \\
$I$: $x<1$, $s_1<2$, $s_2<2$, $s_1+s_2\ge2$ & accept (weight-dependent plans) & reject & dips breach the floors; the reserve cannot finance phasing \\
$x<1$, $s_1+s_2<2$ & reject & reject & even the index fails \\
\bottomrule
\end{tabularx}
\end{table}

\begin{figure}[!htbp]
\centering
\includegraphics[width=\linewidth]{fig_witness_triangle.png}
\caption{The discrepancy triangle $\Q=\{s_1<2,\ s_2<2,\ s_1+s_2\ge2\}$ in the floor plane. Panel A ($x<1$): $\Q$ is the impossibility region $I=\FP$; the arrows show the worst-case dips of \act{fast} ($s_1\to s_1-2$) and \act{slow} ($s_2\to s_2-2$) crossing the floors from the witness point $(6/5,6/5)$, and the point $(1/10,1/10)$ is the second strictness witness. Panel B ($x\ge1$): the same triangle is the rescue set $R$, served by \act{staged}. Below the diagonal both assessments reject when $x<1$ and accept when $x\ge1$. The reserve bars show the bridge stock against the rescue cost $c=1$ at representative slices ($x=0$ and $x=1$; the witness point sits at $x=\tfrac12$). The bottom strip gives the licensing zones at the witness point: \act{slow} only for $r<\rho_1=\tfrac23$, both for $\tfrac23\le r\le\tfrac32$, \act{fast} only for $r>\rho_2=\tfrac32$.}
\label{fig:triangle}
\end{figure}

\begin{figure}[!htbp]
\centering
\includegraphics[width=0.9\linewidth]{fig_path_view.png}
\caption{Path view at the witness point under the worst-case disturbance. Rows: \act{fast}, \act{slow}, \act{staged}. Left: the floor paths (solid) against the zero floor (dashed). Right: the scalarized index $s_1+rs_2$ for $r\in\{2,1,\tfrac12\}$. Under \act{fast} the index stays nonnegative for $r=2$ and $r=1$ while $s_1$ breaches its floor; under \act{slow} the mirror holds. The only plan that protects both floors, \act{staged}, consumes the reserve linearly: infeasible at $x=\tfrac12$, feasible at $x=1$ -- the difference between $I$ and $R$.}
\label{fig:paths}
\end{figure}

\subsection{Persistence under horizon extension, and its limit}\label{sec:persist}

Is the one-interval separation an artefact of the one-interval framing? Two results delimit the answer.

\begin{proposition}[persistence and erasure]\label{prop:persist}
\textup{(i)} Extend the witness datum to $m\ge2$ review intervals by prepending hold intervals, with sole prefix action \act{hold} (tube $\{z\}$, successor $\{z\}$, safe set $S_0$) available at every prefix state, unchanged terminal and successor architecture, and no additional reset or disturbance branches. Then for every stage $j$, $\mathcal V^{\mathrm{typ}}_j\subseteq\bigcap_{w}\mathcal V^{w}_j\subseteq\mathcal V^{\mathrm{phys}}_j$, the stage-0 accepted regions are the witness regions pulled back through the holds, and both strictness witnesses of Theorem~\ref{thm:witness}(5) persist.
\textup{(ii)} Persistence is architecture-conditional: there is a two-interval datum whose single-interval truncation exhibits an acceptance gap and whose two-stage extension erases it, because the second stage has coincident typed and scalarized terminal sets and a repair action admissible from every state.
\end{proposition}

The erasure datum of (ii) (SI~S2.2) has two floors, no reserve, a first stage with actions that move one unit from one floor to the other, and a second stage with a universal repair action; at $(\tfrac25,\tfrac25)$ the truncation has the gap (the two first-stage actions are licensed for $r\ge\tfrac37$ and $r\le\tfrac73$ respectively, and neither alone is typed-safe), while the extension accepts the state under both doctrines. The witness datum itself is not subject to this erasure: its gap is tube-driven -- on $I$ every plan violates the stage-local path constraint $S_0$ -- and a later stage cannot repair a violated tube. Together the two parts state exactly what is claimed: hold-prefix extension preserves the separation; arbitrary multi-stage extension does not; and the erasure mechanisms (a common typed-safe continuation for every aggregate-feasible state, coincident terminal sets on the reachable set, or actions coupled across stages) are realisable.

\subsection{What closes the gap and what does not}\label{sec:closing}

Theorem~\ref{thm:menu} predicts that convexifying the menu closes the gap. On the witness datum the convexification instrument is the \emph{blend family}: for $\delta\in(0,1)$, $\act{blend}_\delta$ is the plan obtained by fractional allocation of the primitive control flows, $u_\delta=\delta\,u_{\act{fast}}+(1-\delta)\,u_{\act{slow}}$; by linearity of the flow map in the control, its worst-case tube is the pointwise convex combination of the two primitive worst-case tubes and its successor the convex combination of the two successors. This is menu convexification at the level of control inputs -- an explicit model assumption, which must be demonstrated from the model equations or carried as a management assumption (fractional allocation of effort, quota or season, with monitoring and enforcement) before the theorem is applied to a real menu.

\begin{theorem}[blend collapse]\label{thm:blend}
On the witness datum, over $X_0$:
\begin{enumerate}[label=\textup{(\roman*)},leftmargin=2em]
\item $\act{blend}_\delta$ is typed-admissible at $z$ exactly when $\delta\in\big[\,1-\tfrac{s_2}2,\ \tfrac{s_1}2\,\big]\cap[0,1]$, an interval nonempty exactly when $s_1+s_2\ge2$ (a singleton at equality), independent of $x\ge0$ and of the weight;
\item the typed accepted set of the menu augmented by the blend family is exactly the compensatory region,
$\Vtyp^{\mathrm{blend}}=\{x\ge1\}\cup\{s_1+s_2\ge2\}=\Vweak$: the gap vanishes, and the blend family reaches nothing beyond $\Vweak$;
\item the admissible window depends on the state through $(s_1,s_2)$ alone, so a single $\delta$ serves every weight: on $I$ the blend is the common plan that no member of the finite menu supplies.
\end{enumerate}
\end{theorem}

\begin{proposition}[discrete time-sharing does not close the gap]\label{prop:timeshare}
Let a sequential policy alternate between \act{fast} and \act{slow} over the interval, each held at full strength on its sub-intervals, so that the visited set is the union of the two primitive worst-case tubes. Such a policy is typed-admissible if and only if $s_1\ge2$ and $s_2\ge2$; on $I$ none is, and the gap survives.
\end{proposition}

The pair delimits ``mixing'' precisely. Continuous convexification -- fractional allocation of flows -- closes the gap exactly at the compensatory region $s_1+s_2\ge2$; discrete alternation at plan granularity closes it only where both dips are separately survivable. The structure is that of pure versus mixed strategies in the assessor--planner game \citep{vonNeumann1928,Sion1958}: with the planner choosing a plan, the assessor a weight, and the payoff the worst-case aggregate, common-plan acceptance is the pure-strategy value, per-weight acceptance the mixed value, and the gap the pure-strategy duality gap that the convex extension closes; its reach--avoid form is the discriminating kernel of \citet{Cardaliaguet1996}. Proposition~\ref{prop:timeshare} is scoped to worst-case alternation at finite rate; infinitesimally fast switching recovers the convexified trajectory in the sense of relaxed controls \citep{Filippov1988,Warga1972} and is subsumed by Theorem~\ref{thm:blend}, not distinct from it.

The remaining instrument acts on the reserve. Define the augmentation map $\mathsf{Aug}_\kappa:(x,s,\mathcal A)\mapsto(x+\kappa,\ s,\ \mathcal A\cup\{\act{staged}_\kappa\})$, where $\act{staged}_\kappa$ is \act{staged} run from the augmented reserve, and the \emph{rescue threshold} $\kappa^*(z)=\inf\{\kappa:\ \exists a\in\mathcal A_\kappa,\ a\in\Etyp(z)\}$.

\begin{proposition}[rescue threshold]\label{prop:rescue}
On the witness datum, for every $z\in X_0$,
\[
\kappa^*(z)=(1-x)_+\cdot\mathbf 1[\,x<1,\ s_1<2,\ s_2<2\,]=(1-x)\cdot\mathbf 1[\,z\notin\Vtyp\,].
\]
\end{proposition}

On $I$ the threshold is $1-x>0$; on $R$ and on all of $\Vtyp$ it is $0$. Because $\act{staged}_\kappa$ is admissible exactly when $x+\kappa\ge1$, independently of the floors, the reserve route admits every state of the failure set $\{x<1,s_1<2,s_2<2\}$ at cost $1-x$ -- including states such as $(\tfrac12,\tfrac1{10},\tfrac1{10})$ that even the aggregate rejects; $I$ is the part of the failure set on which the aggregate still certifies. Along this route the distance to the feasible slice $\{x\ge1\}$ equals the violation $(1-x)_+$ itself, so the error-bound modulus \citep{Fabian2010} of the relation is identically $1$: the route is perfectly conditioned. Table~\ref{tab:instruments} summarises how the instruments act on the gap.

\begin{table}[!htbp]
\centering\small
\caption{Instruments acting on the acceptance gap.}
\label{tab:instruments}
\begin{tabularx}{\textwidth}{@{}>{\raggedright\arraybackslash}p{0.24\textwidth}>{\raggedright\arraybackslash}p{0.32\textwidth}X@{}}
\toprule
Instrument & Operation & Effect on the gap \\
\midrule
Reweighting & change $w$ & none: the gap is a statement over all $w$ \\
Restricting the weight family & replace $\Wp$ by $W\subset\Wp$ & enlarges per-weight acceptance; leaves the common-plan problem (Proposition~\ref{prop:monotone}) \\
Reserve augmentation & add $\act{staged}_\kappa$, $x+\kappa\ge1$ & closes it where only the reserve binds, at cost $\kappa^*=1-x$ (Proposition~\ref{prop:rescue}) \\
Menu convexification & replace $\mathcal D$ by $\conv(\mathcal D)$ (blend family) & closes it exactly on the scalarized region (Theorems~\ref{thm:menu}, \ref{thm:blend}) \\
Discrete time-sharing & union of the primitive tubes & closes it only where both dips are separately survivable (Proposition~\ref{prop:timeshare}) \\
\bottomrule
\end{tabularx}
\end{table}

\subsection{Beyond linear aggregation: the substitutability ladder}\label{sec:ladder}

The results so far concern the linear index. A natural objection is that the gap is an artefact of perfect substitutability and would close under an aggregator with limited substitution. We evaluate Protocol~2 under the CES family. Work in floor-referenced indices $\lambda_i=1+s_i$ ($\lambda_i=1$ at the floor, $\lambda_i\le0$ a full floor-unit below it, the currency of geometric-mean abundance indices) and, for weights normalised to sum to one, the power means
\[
M_\theta(\lambda;w)=\Big(\sum_iw_i\lambda_i^\theta\Big)^{1/\theta},\qquad\sigma=\frac1{1-\theta},
\]
with the usual limits: $\theta=1$ is the linear aggregate, $\theta=0$ the geometric mean (the functional form of the Living Planet Index), $\theta=-1$ the harmonic mean, $\theta\to-\infty$ the Leontief minimum. A state is \emph{accepted at rung $\theta$} when, for every weight, some plan of the menu keeps $M_\theta\ge1$ along its worst-case tube; the physical and destination constraints are aggregator-independent (\act{staged} needs $x\ge1$ and then serves every weight at every rung). Two conventions fix the family. With normalised weights $\sum_iw_i\lambda_i\ge1\iff w\cdot s\ge0$, so the $\theta=1$ member reproduces the compensatory operator of Section~\ref{sec:setting} exactly, with the same licensing thresholds (the \emph{handshake}). For every $\theta<1$ a tube value $\lambda_i\le0$ rejects the plan (the \emph{collapse convention}: the geometric mean is $0$ there and the $\theta<0$ members are undefined there); only the linear member compensates across a collapsed coordinate.

On the gap-region diagonal $z=(x<1,s,s)$, $s\in(1,2)$, the per-weight cover condition reduces, by the symmetry of the two principal plans, to a single comparison at the equal weight: accept iff $(s-1)^\theta+(s+1)^\theta\ge2$ for $\theta>0$, iff $s^2\ge2$ for $\theta=0$, and iff $(s-1)^\theta+(s+1)^\theta\le2$ for $\theta<0$. Each rung therefore has an exact critical floor $s^*(\theta)$, and the critical floors form a ladder (Table~\ref{tab:ladder}, Figure~\ref{fig:ladder}).

\begin{theorem}[the substitutability ladder]\label{thm:ladder}
On the witness datum:
\begin{enumerate}[label=\textup{(\roman*)},leftmargin=2em]
\item \textup{(Nesting.)} For rungs $\theta>\theta'$ and every state, acceptance at $\theta'$ implies acceptance at $\theta$. The accepted rungs at each state form a prefix of the ladder, so the \emph{critical elasticity} $\sigma^*(z)=\inf\{\sigma:\ z\text{ accepted at }\sigma\}$ is well defined; it is strictly positive and finite on every gap state with both floors above $1$, and equal to $\infty$ on gap states with $\min(s_1,s_2)\le1$, which only the linear member certifies.
\item \textup{(Pointwise.)} For every gap state there is a strictly positive elasticity at which the aggregator rejects it: full Leontief is never necessary at a fixed state.
\item \textup{(Uniform.)} For every $\theta<1$ the rung falsely certifies an exact interval $[s^*(\theta),2)$ of diagonal gap states, with $s^*(\theta)<2$ strictly; the critical floors are strictly increasing in depth, $1<5/4<\sqrt2<\varphi<\sqrt3<\cdots\to2$, and converge to the typed boundary. Hence no positive elasticity is uniformly safe on the gap region, $\inf_z\sigma^*(z)=0$, and the only uniformly safe member is $\sigma=0$ -- the Leontief member, which is exactly the typed operator ($\mathcal V^{0}=\Vtyp$).
\end{enumerate}
\end{theorem}

\begin{table}[!htbp]
\centering\small
\caption{The exact ladder on diagonal gap states $z=(x<1,s,s)$: acceptance condition and critical floor $s^*(\theta)$ per rung.}
\label{tab:ladder}
\begin{tabularx}{\textwidth}{@{}ll>{\raggedright\arraybackslash}p{0.4\textwidth}X@{}}
\toprule
$\theta$ & $\sigma$ & accept iff & critical floor $s^*(\theta)$ \\
\midrule
$1$ & $\infty$ & $s\ge1$ & $1$ \\
$2/3$ & $3$ & $s^2\ge3$, else $27(s^2-1)^2\ge(3-s^2)^3$ & root in $(117/100,\,59/50]$ \\
$1/2$ & $2$ & $s\ge5/4$ & $5/4$ \\
$0$ & $1$ & $s^2\ge2$ & $\sqrt2$ \\
$-1$ & $1/2$ & $s^2\ge s+1$ & $\varphi=(1+\sqrt5)/2$ ($8/5$ rejects, $13/8$ accepts) \\
$-2$ & $1/3$ & $s^2\ge3$ & $\sqrt3$ \\
$-m$ & $1/(m+1)$ & $(s-1)^m+(s+1)^m\le2(s^2-1)^m$ & increasing in $m$ \\
$-\infty$ & $0$ & $s\ge2$ & $2$: the typed boundary \\
\bottomrule
\end{tabularx}
\end{table}

\begin{figure}[!htbp]
\centering
\includegraphics[width=\linewidth]{fig_substitutability_ladder.png}
\caption{The substitutability ladder (Theorem~\ref{thm:ladder}). (a) The critical floor $s^*(\theta)$ on the diagonal as a function of the elasticity $\sigma$; the shaded region above the curve is the set of diagonal gap states that the $\sigma$-aggregator falsely certifies. The exact algebraic rungs $1,\,5/4,\,\sqrt2,\,\varphi,\,\sqrt3$ are marked; the curve approaches the typed boundary $s=2$ only as $\sigma\to0$. (b) The critical elasticity $\sigma^*(s)$ along the diagonal (root of the master comparison), with the exact rational brackets of the deposited verifier overlaid: at the witness point $s=6/5$, $\sigma^*\in(2,3)$; at $s=5/4$, $\sigma^*=2$ exactly; $\sigma^*\to\infty$ as $s\to1^+$ and $\sigma^*\to0$ as $s\to2^-$.}
\label{fig:ladder}
\end{figure}

The theorem settles the objection in both directions. Pointwise, the separation is not a Leontief phenomenon: at the witness point every aggregator with $\sigma\le2$ already rejects, and the dashboard that reports $(1/5+11/5)/2=6/5\ge1$ at the trough is contradicted by the geometric index $\sqrt{11/25}<1$ on the same datum. Uniformly, however, limited substitutability does not restore soundness: for every positive elasticity there are gap states, arbitrarily close to the typed boundary, that the aggregator certifies and no plan completes safely. The appearance of $\sqrt2$, $\varphi$ and $\sqrt3$ as the exact critical floors of the geometric, harmonic and $\sigma=1/3$ members is, to our knowledge, a new exact witness for the substitutability debate. Off the diagonal the geometric member's acceptance region is $\{s_2\ge\rho(s_1)\}$ for an explicit decreasing involution $\rho$ with fixed point $\sqrt2$, and the complete two-dimensional landscape of $\sigma^*$ -- including the fact that, at a fixed total margin below $2\sqrt2$, a moderately concentrated deficit is \emph{more} certifiable than a balanced one -- is decided by certified rational enclosures in SI~S3. The identification of the $\theta=0$ member with the Living Planet Index is at the level of functional form only; the index's published inputs and weighting are not Protocol~2, and the rung decision is the adversarial version of its reading (SI~S3.5).

\subsection{Beyond action-indexed disturbances: the common shock}\label{sec:commonshock}

A second objection is that the gap depends on the action-indexed disturbance convention, under which the adverse event strikes only the coordinate the active plan moves. Decompose the datum as on the benchmark: each principal plan's announced schedule is a mid-interval tent of depth $3/2$ in its characteristic coordinate, and each environmental event is a transient tent dip of depth $\delta$ in one floor coordinate (the datum's depth $2=3/2+1/2$ fixes $\delta=1/2$); \act{staged}'s floors grow linearly, so an event of depth $\delta$ dips a growing floor by $\max(0,\delta-1/8)$. Under the \emph{common-shock class} the two coordinate events are separate disturbances each of which strikes every plan -- a marine heatwave striking the same biomass coordinate regardless of which harvest plan is running.

\begin{proposition}[the common-shock verdict]\label{prop:commonshock}
Under the common-shock class at $\delta=1/2$:
\begin{enumerate}[label=\textup{(\roman*)},leftmargin=2em]
\item $\Vtyp^{\mathrm{cs}}=\{s_1\ge2,\ s_2\ge\tfrac12\}\cup\{s_2\ge2,\ s_1\ge\tfrac12\}\cup\{x\ge1,\ s_1\ge\tfrac38,\ s_2\ge\tfrac38\}$ and
$\Vweak^{\mathrm{cs}}=\{x<1,\ s_1\ge\tfrac12,\ s_2\ge\tfrac12,\ s_1+s_2\ge2\}\cup\{x\ge1,\ s_1\ge\tfrac38,\ s_2\ge\tfrac38\}$, so the gap
$\mathrm{FP}_{\mathrm{cs}}=\{x<1,\ \tfrac12\le s_1<2,\ \tfrac12\le s_2<2,\ s_1+s_2\ge2\}$
survives: it is the action-indexed gap $I$ reduced by the exact $1/2$-cut, nonempty with interior, and the witness point is a member. On the surviving region the licensing thresholds are unshifted.
\item The rescue threshold becomes a trichotomy: $\kappa^*_{\mathrm{cs}}=0$ on $\Vtyp^{\mathrm{cs}}$; $(1-x)_+$ on failure states with $\min(s_1,s_2)\ge\tfrac38$; and $+\infty$ on failure states with $\min(s_1,s_2)<\tfrac38$, which no reserve increment rescues (witness $(\tfrac12,\tfrac14,3)$).
\item The blend window is nonempty exactly on $\{s_1\ge\tfrac12,\ s_2\ge\tfrac12,\ s_1+s_2\ge\tfrac52\}$, so the common-shock gap is not blend-closed on the band $s_1+s_2\in[2,\tfrac52)$ (witness $(\tfrac12,1,1)$).
\end{enumerate}
\end{proposition}

The separation is therefore not an artefact of the action-indexed convention; what the shared shock changes is the remedy. On the $1/2$-cut the compensatory verdict itself flips from accept to reject, and on $\{\min(s)<3/8\}$ the prescription changes from accumulating reserve to rebuilding the weaker floor first, because no reserve increment certifies. The coupled class (one disturbance carrying both events), its depth dependence with exact boundary $\delta=5/4$, the relocation of the gap to larger floors at full magnitude, and the collapse of the separation under a replace reading in which every plan faces the same plan- and weight-independent dip are stated and proved in SI~S4; the mechanism of that collapse is instructive -- when every plan faces the same dip, weight-wise plan choice buys nothing, and the per-weight cover degenerates to the typed condition.

\subsection{Verification}\label{sec:verification}

Every continuum statement above is proved: Lemmas~\ref{lem:quant}--\ref{lem:cone}, Propositions~\ref{prop:hierarchy}--\ref{prop:monotone} and Theorem~\ref{thm:menu} in \ref{app:general}; Theorem~\ref{thm:witness}, Theorem~\ref{thm:blend} and Propositions~\ref{prop:timeshare}--\ref{prop:rescue} in \ref{app:witness}; Proposition~\ref{prop:persist}, Theorem~\ref{thm:ladder} and Proposition~\ref{prop:commonshock} in SI~S2. Three further layers check the finite rational content. (a) A deterministic exact-integer grid verifier (scale 40; no floating point, tolerances or randomness) checks the plan classifications, the region identities and the accepted-set identities of Theorem~\ref{thm:witness} and Proposition~\ref{prop:persist}(i) on a $31^3=29{,}791$-state grid over $(x,s_1,s_2)$, with the critical ratios $\rho_1,\rho_2$ and their midpoint tested at every grid state; its 25 checks are enumerated in SI~S7. (b) Separate fail-loud exact-rational verifiers (standard-library fractions only, byte-reproducible logs) check the ladder and its off-diagonal closure (57, 50 and 44 checks), the common-shock variant (77 checks, with a handshake that reproduces Theorem~\ref{thm:witness} under the action-indexed class on 676 grid states) and the benchmark of Section~\ref{sec:benchmark} (24 checks); the lists are not pooled. (c) The theorem layer is formalised in Lean~4: the witness datum exactly, Theorem~\ref{thm:witness}(1)--(7), Lemma~\ref{lem:cone} and Proposition~\ref{prop:hierarchy} on the datum, Theorem~\ref{thm:blend} and Proposition~\ref{prop:timeshare}, the handshake and collapse conventions, the master-equation reductions and Theorem~\ref{thm:ladder}(ii)--(iii) with its false-certification witnesses, in a module of 199 theorems that compiles with no \texttt{sorry}, no additional axioms and no computer-algebra oracle, at an ordered-field interface of which the real-number statements are an instance. The Lean layer and the appendix proofs were produced independently and agree; neither replaces the other, and the grid check is a third and separate thing again, since it validates the symbolic classification on one enumerated instance, which is all a finite check can do. Code, logs and figure pipelines are deposited \citep{Abaee2026code}.

